\makeatletter
\def\input@path{{preprints-template/preprints_template_latex/}}
\let\preprint@bibliographystyle\bibliographystyle
\renewcommand{\bibliographystyle}[1]{%
  \preprint@bibliographystyle{preprints-template/preprints_template_latex/#1}}
\makeatother

\documentclass[preprints,article,submit,oneauthor]{Definitions/mdpi}
\setlength{\headheight}{18pt}

\makeatletter
\let\bibliographystyle\preprint@bibliographystyle
\makeatother

\graphicspath{{preprints-template/preprints_template_latex/}}

\firstpage{1}
\makeatletter
\setcounter{page}{\@firstpage}
\makeatother
\pubvolume{1}
\issuenum{1}
\articlenumber{0}
\pubyear{2026}
\copyrightyear{2026}
\datereceived{}
\daterevised{}
\dateaccepted{}
\datepublished{}

\newcounter{manuscriptresult}[section]
\renewcommand{\themanuscriptresult}{\thesection.\arabic{manuscriptresult}}
\theoremstyle{mdpi}
\newtheorem{theorem}[manuscriptresult]{Theorem}
\newtheorem{lemma}[manuscriptresult]{Lemma}
\newtheorem{proposition}[manuscriptresult]{Proposition}
\newtheorem{corollary}[manuscriptresult]{Corollary}

\theoremstyle{definition}
\newtheorem{remark}[manuscriptresult]{Remark}

\newcommand{\C}{\mathbb C}
\newcommand{\D}{\mathbb D}
\newcommand{\Mn}{\mathbb M_n(\C)}
\newcommand{\alg}{\operatorname{alg}}
\newcommand{\distop}{\operatorname{dist}}
\newcommand{\RePart}{\operatorname{Re}}
\newcommand{\diag}{\operatorname{diag}}
\newcommand{\iu}{\mathrm{i}}

\Title{The numerical range is a 2-spectral set}
\Author{Shanmu Jin}
\AuthorNames{Shanmu Jin}
\address{Department of Neurosurgery, Peking Union Medical College
  Hospital, No.~1 Shuaifuyuan, Dongcheng District, Beijing 100730,
  China\\
jinshanmu1129@pku.edu.cn}

\abstract{\texorpdfstring{%
For an $n\times n$ complex matrix $A$, let
$W(A)=\{x^*Ax:x\in\C^n,\ \lVert x\rVert_2=1\}$ be its numerical range.
Here $x^*$ denotes the conjugate transpose, $\lVert\,\cdot\,\rVert_2$
the Euclidean norm, and $\lVert\,\cdot\,\rVert$ the induced operator
norm.  For every complex polynomial $p$, we prove that
\[
  \lVert p(A)\rVert\leq
  2\max_{z\in W(A)}\lvert p(z)\rvert.
\]
More generally, the same estimate holds for every function $f$
holomorphic on a neighborhood of $W(A)$, with $p$ replaced by $f$;
consequently, $W(A)$ is a $2$-spectral set for $A$.  This
proves Crouzeix's conjecture, and the constant is optimal.  The central
ingredient is a mass-parameterized positive-real completion theorem
relative to an auxiliary eigenbasis: if a normalized matrix-valued
Carath\'eodory function completes
$\frac{2}{\mathfrak m}(I-wT)^{-1}$ modulo the adjoint algebra, then
$\lVert T\rVert\leq\mathfrak m$ for $\mathfrak m\geq2$.  The
numerical-range double layer has mass two.
The classical positive double-layer calculus supplies such completions
for $f(B)$ whenever the auxiliary matrix $B$ has simple spectrum; the
eigenvalues of $f(B)$ may repeat.  Sampling the associated Herglotz
kernel at scaled conjugate diagonal entries and at the origin cancels
the nonconstant completion term and leaves two ordered weighted
Gramians; their first nonconstant terms yield the mass bound.
Simple-spectrum
approximation and canonical convex outer domains with supported oriented
radial boundaries yield the general theorem.  The mass formulation also
provides a concrete scalar proof protocol for annular, multiply
connected, operator-radius, and abstract spectral-constant problems.
}{%
For a complex square matrix A, we prove the sharp polynomial
Crouzeix inequality and its extension to every function f holomorphic
near the numerical range of A.  Consequently, the numerical range is
a 2-spectral set, and the factor 2 is optimal.  The proof uses a
mass-parameterized positive-real
completion theorem relative to an auxiliary eigenbasis.  A
matrix-valued Herglotz kernel is sampled so that the adjoint-algebra
correction cancels, leaving ordered weighted Gramians that yield the
sharp mass bound.  Simple-spectrum approximation and convex outer
domains give the general result.  The mass formulation also supplies
a reusable scalar protocol for sharp annular, operator-radius, and
abstract spectral-constant problems.
}}

\keyword{Crouzeix conjecture; numerical range; spectral set;
matrix-valued Herglotz kernel; double-layer potential}
\MSC{Primary 47A25; Secondary 47A12, 15A60}

\begin{document}

\section{Introduction}
\label{sec:introduction}

Let $\C$ denote the complex field and, for $n\geq1$, let $\Mn$ denote
the algebra of $n\times n$ complex matrices.  For $A\in\Mn$, its
numerical range is
\begin{equation}
  W(A)=\{x^*Ax:x\in\C^n,\ \lVert x\rVert_2=1\}.
  \label{eq:numerical-range}
\end{equation}
Here $x^*$ denotes the conjugate transpose, $\sigma(A)$ denotes the
spectrum, $\lVert x\rVert_2$ is the Euclidean norm, $\lVert X\rVert$ is
the induced operator norm, and scalar moduli are denoted by
$\lvert\,\cdot\,\rvert$.  The set $W(A)$ is compact, the
Toeplitz--Hausdorff theorem makes it convex, and
$\sigma(A)\subset W(A)$.  Functions holomorphic near $\sigma(A)$ are
evaluated using the holomorphic functional calculus.

A compact set $K\subset\C$ containing $\sigma(A)$ is called a
$C$-spectral set for $A$ if
\begin{equation}
  \lVert r(A)\rVert\leq C\max_{z\in K}\lvert r(z)\rvert
  \label{eq:spectral-set-definition}
\end{equation}
for every rational function $r$ pole-free on $K$.  The following
finite-dimensional theorem states the customary polynomial form of
Crouzeix's conjecture together with its holomorphic and spectral-set
extensions.

\begin{theorem}[Finite-dimensional form]
\label{thm:main}
Let $n\geq1$ and let $A\in\Mn$.  For every $p\in\C[z]$,
\begin{equation}
  \lVert p(A)\rVert
  \leq2\max_{z\in W(A)}\lvert p(z)\rvert.
  \label{eq:matrix-polynomial-bound}
\end{equation}
More generally, if $f$ is holomorphic on a neighborhood of $W(A)$,
then
\begin{equation}
  \lVert f(A)\rVert
  \leq2\max_{z\in W(A)}\lvert f(z)\rvert.
  \label{eq:main-bound}
\end{equation}
Consequently, $W(A)$ is a $2$-spectral set for $A$.
\end{theorem}

The factor $2$ is optimal: for
\begin{equation}
  J=\begin{pmatrix}0&1\\0&0\end{pmatrix}
  \quad\text{and}\quad p(z)=z,
  \label{eq:sharp-example}
\end{equation}
one has
$\max_{z\in W(J)}\lvert z\rvert=1/2$ and $\lVert J\rVert=1$.

For later use, write
\[
  \D=\{z\in\C:\lvert z\rvert<1\}.
\]
We use $I$ for the identity matrix of the size dictated by context,
$X^*$ for the adjoint, put $\RePart X=(X+X^*)/2$, and let $\alg(T)$
denote the unital algebra generated by $T$.  For Hermitian matrices
$X$ and $Y$, we write $X\preceq Y$ (equivalently, $Y\succeq X$) when
$Y-X$ is positive semidefinite.

Crouzeix proved the estimate with constant $2$ for $2\times2$ matrices and
formulated the general conjecture in~\cite{Crouzeix2004}.  His
subsequent dimension-free theorem gave the constant $11.08$~%
\cite{Crouzeix2007}.  After further structural and numerical work~%
\cite{Crouzeix2016}, Crouzeix and Palencia proved that $W(A)$ is a
$(1+\sqrt2)$-spectral set~\cite{CrouzeixPalencia2017}.  Their
double-layer and Cauchy-transform argument was clarified in~%
\cite{RansfordSchwenninger2018} and developed in several directions in~%
\cite{CaldwellGreenbaumLi2018,CrouzeixGreenbaum2019}.

The intervening literature established the conjectured constant for
perturbed Jordan blocks and, more generally, scalar translates of
cyclic weighted shifts, as well as for tridiagonal $3\times3$ matrices
with elliptic numerical range centered at an eigenvalue~%
\cite{ChoiGreenbaum2012,Choi2013,GladerKurulaLindstrom2018}.
Numerical and extremal investigations appear in~%
\cite{GreenbaumOverton2018,Li2020}, and a broad account of the problem
and its variants is given in~\cite{BickelEtAl2020}.  Subsequent work
has treated Blaschke-product level sets, cyclicity reductions,
compressions of shifts, further nilpotent classes, and quantitative
bounds for KMS matrices~%
\cite{BickelGorkin2024,OLoughlinVirtanen2024,BickelEtAl2024,
CrouzeixGreenbaumLi2024}.  An alternative proof for a class of weighted shifts
appears in~\cite{CrouzeixGreenbaum2025}.  Across these cited
reductions, the best dimension-free universal bound remained
$1+\sqrt2$.

\subsection{The symmetrized calculus}

The positive boundary measure used below originates in the
convex-domain integral calculus of Delyon--Delyon~%
\cite{DelyonDelyon1999}; see also~%
\cite{BadeaCrouzeixDelyon2006,CrouzeixPalencia2017}.  If $B$ has
numerical range in a convex domain $\Omega$ with continuously
differentiable boundary and $f$ is holomorphic on a neighborhood of
$\overline\Omega$ with $\lvert f\rvert\leq1$ on $\overline\Omega$, this
calculus produces a unital positive map $\Phi$ and a holomorphic Cauchy
companion $g_f$ such that
\begin{equation}
  2\Phi(f)=f(B)+g_f(B)^*.
  \label{eq:intro-symmetrized}
\end{equation}
Thus positivity directly controls a symmetrized functional calculus,
while the target estimate concerns $f(B)$.  The
Crouzeix--Palencia argument couples the two terms
in~\eqref{eq:intro-symmetrized}; the example
in~\cite{RansfordSchwenninger2018} shows that its abstract norm estimate
is sharp at $1+\sqrt2$.  The corresponding general uniform-algebra
theorem has the same sharp constant; its additional unital case was
proposed as a possible route to $2$
in~\cite{ClouatreOstermannRansford2023}.
Schwenninger and de Vries revisited the double-layer method~%
\cite{SchwenningerDeVries2025}.  Malman et al.\ proved that a
configuration-constant refinement gives a strict improvement for every
fixed numerical-range shape; they also constructed thin quadrilaterals
for which the resulting bounds approach
$1+\sqrt2$~\cite{MalmanEtAl2025}.
Separately,
continuity and compactness give a constant $C_n<1+\sqrt2$ in every
fixed matrix dimension $n$~\cite{MalmanRatio2024}.

The additional information used here is algebraic.  For the approximants
constructed in Section~\ref{sec:passage}, the companion
in~\eqref{eq:intro-symmetrized} belongs to the algebra generated by
$B^*$, where the simple-spectrum matrix $B$ supplies an eigenbasis that
also diagonalizes $f(B)$.  Treating $f(B)$ and the companion jointly
preserves this coupling.

\subsection{Auxiliary-basis de-symmetrization}

Our numerical-range argument uses the mass-two specialization of the
general completion theorem proved below and the whole Cayley family
associated with~\eqref{eq:intro-symmetrized}.  It has four steps.

\begin{enumerate}
\item Applying $\Phi$ to
\[
  \zeta\longmapsto\frac{1+w f(\zeta)}{1-w f(\zeta)},
  \qquad w\in\D,
\]
produces a matrix-valued Carath\'eodory function $H$, meaning that
$H(0)=I$ and $\RePart H(w)\succeq0$.

\item For a simple-spectrum auxiliary matrix $B$, put $T=f(B)$.  The
double-layer identity gives the positive-real completion
\begin{equation}
  H(w)-(I-wT)^{-1}\in\alg(B^*),\qquad w\in\D.
  \label{eq:intro-completion}
\end{equation}

\item The eigenbasis of $B$ diagonalizes $T=f(B)$ even when values of
$f$ on $\sigma(B)$ coincide, and it turns the adjoint-algebra defect
in~\eqref{eq:intro-completion} into a diagonal analytic correction.
Sampling the Herglotz kernel at half the conjugate diagonal entries of
$T$, together with one vector sample at the origin, cancels the
correction exactly; kernel positivity accommodates repeated sampling
points.

\item The surviving positive-matrix inequality compares two weighted
Gramians at scales $\tau$ and $\tau^2$, where
$\tau=\mathfrak m^{-1}$.  Their ordered difference and the first
nonconstant Gramian term yield $\lVert T\rVert\leq\mathfrak m$; here
$\mathfrak m=2$.
\end{enumerate}

The proof uses two limits: simple-spectrum matrices tend to $A$,
followed by outer domains converging to $W(A)$.  Together they yield
the general case.

\section{Positive-real kernels}
\label{sec:herglotz}

Let $\mathcal Q$ be an $\Mn$-valued kernel on $\D\times\D$.  We call
$\mathcal Q$ positive if, for every finite choice
$w_0,\ldots,w_m\in\D$ and $\xi_0,\ldots,\xi_m\in\C^n$,
\begin{equation}
  \sum_{i,j=0}^m
  \xi_i^*\mathcal Q(w_i,w_j)\xi_j\geq0.
  \label{eq:positive-kernel-definition}
\end{equation}

We shall use the following standard matrix-valued Herglotz kernel
fact.  A direct proof is included to fix the normalization.

\begin{lemma}[Matrix Herglotz kernel]
\label{lem:herglotz-kernel}
Let $F\colon\D\to\Mn$ be analytic with
$\RePart F(w)\succeq0$ for every $w\in\D$.  Then
\begin{equation}
  \mathcal L_F(w,z)
  =\frac{F(w)+F(z)^*}{1-w\overline z}
  \label{eq:herglotz-kernel}
\end{equation}
is a positive kernel.
\end{lemma}

\begin{proof}
Fix $w_0,\ldots,w_m\in\D$ and
$\xi_0,\ldots,\xi_m\in\C^n$, and put
\[
  \eta(\zeta)=\sum_{j=0}^m
  \frac{\xi_j}{1-\overline{w_j}\zeta},
  \qquad \lvert\zeta\rvert=1.
\]
For $0\leq r<1$, the scalar Cauchy formula, applied entrywise, gives
\begin{equation}
  \frac{1}{2\pi}
  \int_0^{2\pi}
  \eta(e^{\iu t})^*\,2\RePart
  F(re^{\iu t})\,
  \eta(e^{\iu t})\,dt
  =
  \sum_{i,j=0}^m \xi_i^*
  \frac{F(rw_i)+F(rw_j)^*}
       {1-w_i\overline{w_j}}\xi_j.
  \label{eq:regularized-herglotz-average}
\end{equation}
The integrand on the left is nonnegative, so the right side is
nonnegative.  Letting $r\uparrow1$ and using the continuity of $F$ at
the finitely many points $w_i$ gives
\eqref{eq:positive-kernel-definition} for
$\mathcal Q=\mathcal L_F$.
\end{proof}

\section{The positive-real completion principle}
\label{sec:completion}

The following mass-parameterized completion theorem is the core of the
argument.  It converts a positive-real completion constrained by the
adjoint algebra of an auxiliary simple-spectrum matrix into a norm
bound equal to the mass parameter.

\begin{theorem}[Positive-real completion of mass $\mathfrak m$]
\label{thm:completion}
Let $2\leq\mathfrak m<\infty$, and let
\begin{equation}
  B=S\diag(\beta_1,\ldots,\beta_n)S^{-1},
  \qquad
  T=S\Lambda S^{-1},
  \qquad
  \Lambda=\diag(\lambda_1,\ldots,\lambda_n),
  \label{eq:T-diagonalization}
\end{equation}
where $S$ is invertible, the numbers $\beta_1,\ldots,\beta_n$ are
distinct, and $\lvert\lambda_i\rvert\leq1$ for every $i$.  The numbers
$\lambda_i$ may repeat.
Let $H\colon\D\to\Mn$ be analytic and satisfy
\begin{equation}
  H(0)=I,\qquad \RePart H(w)\succeq0\quad(w\in\D),
  \label{eq:completion-positive-real}
\end{equation}
and
\begin{equation}
  H(w)-\frac{2}{\mathfrak m}(I-wT)^{-1}
  \in\alg(B^*)\qquad(w\in\D).
  \label{eq:completion-algebra-condition}
\end{equation}
Then $\lVert T\rVert\leq\mathfrak m$.
\end{theorem}

\begin{proof}
Put $\tau=\mathfrak m^{-1}$, so $0<\tau\leq1/2$, and set
$G=S^*S\succ0$.
Since the $\beta_i$ are distinct, polynomial interpolation shows that
\begin{equation}
  \alg(B^*)
  =\left\{(S^{-1})^* D S^* \,:\,
  D\ \text{is diagonal}\right\}.
  \label{eq:adjoint-algebra-diagonal}
\end{equation}
By~\eqref{eq:adjoint-algebra-diagonal}, the matrix
\[
  \Theta(w)=S^*\bigl(H(w)-2\tau(I-wT)^{-1}\bigr)(S^*)^{-1}
\]
is diagonal for every $w$.  On setting
$\widetilde H(w)=S^*H(w)S$, we obtain
\begin{equation}
  \widetilde H(w)=2\tau G(I-w\Lambda)^{-1}+\Theta(w)G.
  \label{eq:Htilde-decomposition}
\end{equation}
Since $\widetilde H(0)=G$, one has $\Theta(0)=(1-2\tau)I$.  Write
\[
  \Theta(w)=(1-2\tau)I+\Psi(w),
  \qquad
  \Psi(w)=\diag(\psi_1(w),\ldots,\psi_n(w)),
\]
where $\Psi(0)=0$.  Thus
\begin{equation}
  \widetilde H(w)=F(w)+\Psi(w)G,
  \qquad
  F(w)=(1-2\tau)G+2\tau G(I-w\Lambda)^{-1}.
  \label{eq:mass-H-decomposition}
\end{equation}
Since $\RePart\widetilde H(w)=S^*(\RePart H(w))S$, the real part of
$\widetilde H(w)$ is positive semidefinite.  Consequently,
Lemma~\ref{lem:herglotz-kernel} applies to $\widetilde H$.

Define Hermitian matrices $P,Q,Y\in\Mn$ by
\begin{equation}
  P_{ij}=\frac{G_{ij}}
  {1-\tau^2\overline{\lambda_i}\lambda_j},
  \qquad
  Q_{ij}=\frac{G_{ij}}
  {1-\tau\overline{\lambda_i}\lambda_j},
  \qquad
  Y=Q-P.
  \label{eq:PQY-definition}
\end{equation}
Here
$\lvert\tau^2\overline{\lambda_i}\lambda_j\rvert\leq1/4$ and
$\lvert\tau\overline{\lambda_i}\lambda_j\rvert\leq1/2$, so the
denominators are nonzero.  Since $G=G^*$, conjugate symmetry in the
displayed formulas gives $P=P^*$ and $Q=Q^*$.
Let $e_1,\ldots,e_n$ be the standard basis of $\C^n$, and fix
$u=(u_1,\ldots,u_n)^T\in\C^n$.  In the positive-kernel inequality use
the $n$ points and vectors
\begin{equation}
  w_i=\tau\overline{\lambda_i},\qquad
  \xi_i=u_i e_i\qquad(1\leq i\leq n),
  \label{eq:eigenvalue-samples}
\end{equation}
and add the point $w_0=0$ with vector
\begin{equation}
  \xi_0=v=-G^{-1}Pu.
  \label{eq:compensating-vector}
\end{equation}
The indexed sample remains valid when points coincide, including
indices for which $\lambda_i=0$.  Its contribution from $\Psi$ to the
sampled kernel quadratic form is
\begin{equation}
  2\RePart\sum_{i=1}^n
  \overline{u_i}\psi_i(w_i)
  \left(
    (Gv)_i+
    \sum_{j=1}^n
    \frac{G_{ij}u_j}{1-w_i\overline{w_j}}
  \right).
  \label{eq:theta-contribution}
\end{equation}
Because
\[
  1-w_i\overline{w_j}
  =1-\tau^2\overline{\lambda_i}\lambda_j,
\]
the parenthesis in~\eqref{eq:theta-contribution} is
$(Gv+Pu)_i$, which vanishes by~%
\eqref{eq:compensating-vector}.

It remains to calculate the contribution of $F$.  For
$x=\overline{\lambda_r}\lambda_s$, the sample--sample entry is
\begin{equation}
  2(1-2\tau)P_{rs}
  +\frac{4\tau G_{rs}}{(1-\tau x)(1-\tau^2x)}.
  \label{eq:sampled-main-block}
\end{equation}
The partial-fraction identity
\[
  \frac1{(1-\tau x)(1-\tau^2x)}
  =\frac1{1-\tau}\left(\frac1{1-\tau x}
  -\frac{\tau}{1-\tau^2x}\right)
\]
shows that the compressed sample--sample block is
\[
  M=2P+\frac{4\tau}{1-\tau}Y.
\]
The sample--origin block is $R_0=2(1-\tau)G+2\tau Q$, and the
origin--origin block is $2G$.  After cancellation of the $\Psi$ term,
kernel positivity on the graph $v=-G^{-1}Pu$ gives
\begin{equation}
  0\preceq M-R_0G^{-1}P-PG^{-1}R_0+2PG^{-1}P.
\label{eq:pre-gramian-inequality}
\end{equation}
Using $Q=P+Y$ reduces this to
\begin{equation}
  \frac{4\tau}{1-\tau}Y
  -2\tau(YG^{-1}P+PG^{-1}Y)
  +2(1-2\tau)(PG^{-1}P-P)\succeq0.
  \label{eq:Y-inequality}
\end{equation}

Balance the diagonalization by defining
\begin{equation}
\begin{aligned}
  \widetilde T&=G^{1/2}\Lambda G^{-1/2},
  &\widehat P&=G^{-1/2}PG^{-1/2},\\
  \widehat Q&=G^{-1/2}QG^{-1/2},
  &\widehat Y&=\widehat Q-\widehat P.
\end{aligned}
\label{eq:balanced-definitions}
\end{equation}
Expanding the scalar kernels in~\eqref{eq:PQY-definition} into
geometric series gives
\begin{equation}
  \widehat P=\sum_{k=0}^{\infty}
  \tau^{2k}\widetilde T^{*k}\widetilde T^k,
  \qquad
  \widehat Q=\sum_{k=0}^{\infty}
  \tau^k\widetilde T^{*k}\widetilde T^k.
  \label{eq:two-gramians}
\end{equation}
Indeed, $\widetilde T^k=G^{1/2}\Lambda^k G^{-1/2}$, so the sequence
$(\widetilde T^k)$ is bounded even when some
$\lvert\lambda_j\rvert=1$; the
geometric weights give norm convergence.  It follows that
\begin{equation}
  \widehat P\succeq I,
  \qquad
  \widehat Y=
  \sum_{k=1}^{\infty}(\tau^k-\tau^{2k})
  \widetilde T^{*k}\widetilde T^k
  \succeq\frac{1-\tau}{\tau}(\widehat P-I).
  \label{eq:Yhat-positive}
\end{equation}
The last inequality follows termwise.  Taking the congruence
of~\eqref{eq:Y-inequality} by $G^{-1/2}$ and dividing by $2\tau$ yields
\begin{equation}
  \frac{2}{1-\tau}\widehat Y
  -\widehat Y\widehat P-\widehat P\widehat Y
  +\frac{1-2\tau}{\tau}(\widehat P^2-\widehat P)\succeq0.
  \label{eq:Yhat-inequality}
\end{equation}

Let $e$ be a unit eigenvector of $\widehat P$, say
$\widehat Pe=\alpha e$, and put $y=e^*\widehat Ye$.  Then
$\alpha\geq1$ and~\eqref{eq:Yhat-inequality} gives
\begin{equation}
  0\leq2\left(\frac1{1-\tau}-\alpha\right)y
  +\frac{1-2\tau}{\tau}\alpha(\alpha-1).
  \label{eq:eigenvector-evaluation}
\end{equation}
If $\alpha\leq(1-\tau)^{-1}$, then $\alpha\leq2$.  Otherwise the
coefficient of $y$ is negative, while
$y\geq(1-\tau)(\alpha-1)/\tau$ by~\eqref{eq:Yhat-positive}.  Because the
coefficient is negative, this lower bound for $y$ gives an upper bound
for the right side of~\eqref{eq:eigenvector-evaluation}, namely
\[
  0\leq\frac{\alpha-1}{\tau}(2-\alpha).
\]
Thus $\alpha\leq2$ in either case.  Therefore
\begin{equation}
  I\preceq\widehat P\preceq2I.
  \label{eq:Phat-bounds}
\end{equation}
The $k=1$ term of the same series now gives
\[
  \tau^2\widetilde T^*\widetilde T
  \preceq\widehat P-I\preceq I.
\]
Thus $\lVert\widetilde T\rVert\leq\tau^{-1}=\mathfrak m$.  The polar
decomposition
$S=UG^{1/2}$ gives
\[
  T=U\widetilde TU^*,
\]
where $U$ is unitary.  Therefore
$\lVert T\rVert\leq\mathfrak m$.
\end{proof}

\begin{remark}
\label{rem:completion-novelty}
The sampling scale $\tau=\mathfrak m^{-1}$ simultaneously produces the
weights $\tau$ and $\tau^2$ in~\eqref{eq:two-gramians}.  Their ordered
difference bounds $\widehat P$, and the first nonconstant term then
gives the coefficient $\mathfrak m$.  The origin sample cancels the
nonconstant diagonal correction exactly.  When $\mathfrak m=2$, the
last term of~\eqref{eq:Yhat-inequality} vanishes and the theorem
reduces to the sharp completion used for numerical ranges.  For
$\mathfrak m>2$, the positive constant part
$\Theta(0)=(1-2/\mathfrak m)I$ supplies the additional Gramian term
in~\eqref{eq:Yhat-inequality}.
Distinct auxiliary eigenvalues identify $\alg(B^*)$ with the diagonal
algebra in the chosen basis, while the diagonal entries of $T$ may
repeat.
\end{remark}

\section{The double-layer completion}
\label{sec:double-layer}

We now construct from the numerical range the mass-two function
required by Theorem~\ref{thm:completion}.  The measure and companion
transform used here are classical~%
\cite{DelyonDelyon1999,Crouzeix2007,CrouzeixPalencia2017}.  The full
Cayley family yields the algebraic
constraint~\eqref{eq:completion-algebra-condition} relative to the
auxiliary matrix.

Section~\ref{sec:passage} uses the following canonical parallel
domains.  We say that a bounded convex domain $\Omega$ has a
\emph{supported oriented radial boundary} if, for some $c\in\Omega$,
its boundary is traced once by
\[
  \gamma(t)=c+\rho(t)e^{\iu t},\qquad 0\leq t\leq2\pi,
\]
where $\rho$ is positive, continuously differentiable, and
$2\pi$-periodic, and there are continuous $2\pi$-periodic functions
$\nu$ and $s$ such that $\lvert\nu(t)\rvert=1$, $s(t)\geq0$,
\begin{equation}
  \gamma'(t)=\iu\nu(t)s(t),
  \qquad
  \RePart\bigl(\overline{\nu(t)}(\gamma(t)-z)\bigr)\geq0
  \quad(z\in\Omega).
  \label{eq:oriented-radial-boundary}
\end{equation}
Thus $\nu$ is the outward unit normal, $s(t)\,dt$ is arclength, and
the orientation is positive.  Lemma~\ref{lem:smooth-outer-approximation}
below supplies exactly this boundary for every canonical parallel
domain.

\begin{lemma}[Radial double-layer positive map]
\label{lem:double-layer-map}
Let $B\in\Mn$, let $\Omega\subset\C$ be a bounded convex domain with a
supported oriented radial boundary $\gamma$, and suppose
$W(B)\subset\Omega$.  There is a bounded complex-linear, unital,
positive map on the compact parameter interval,
\[
  \Phi_\gamma\colon C([0,2\pi])\longrightarrow\Mn,
\]
given by the double-layer density below.
\end{lemma}

\begin{proof}
Use the data in~\eqref{eq:oriented-radial-boundary}, and put
\[
  Z_t=\gamma(t)I-B,
  \qquad
  \mathcal R_t=Z_t^{-1}.
\]
This resolvent exists because
$\sigma(B)\subset W(B)\subset\Omega$.  Define the matrix-valued
boundary density
\begin{equation}
  D_B(t)=\frac{s(t)}{2\pi}
  \left(
    \nu(t)\mathcal R_t+
    \overline{\nu(t)}\mathcal R_t^*
  \right).
  \label{eq:double-layer-density}
\end{equation}
It is positive semidefinite.  The support inequality in
\eqref{eq:oriented-radial-boundary} gives
\begin{equation}
  \nu(t)Z_t^*+\overline{\nu(t)}Z_t
  =2\RePart\bigl(\overline{\nu(t)}Z_t\bigr)\succeq0.
  \label{eq:supporting-half-plane}
\end{equation}
Indeed, for every unit vector $x\in\C^n$, the quadratic form on the
left is
\[
  2\RePart\bigl(\overline{\nu(t)}
  (\gamma(t)-x^*Bx)\bigr)\geq0,
\]
because $x^*Bx\in W(B)\subset\Omega$.
Congruence by $\mathcal R_t$ turns the left side of
\eqref{eq:supporting-half-plane} into the matrix in parentheses in
\eqref{eq:double-layer-density}.

Since $d\gamma=\iu\nu(t)s(t)\,dt$, the matrix Cauchy formula gives
\begin{equation}
  \frac{1}{2\pi}\int_0^{2\pi}
  \nu(t)\mathcal R_t s(t)\,dt=I.
  \label{eq:first-layer-mass}
\end{equation}
Taking adjoints gives the same identity for the second layer.  Hence
\begin{equation}
  \int_0^{2\pi}D_B(t)\,dt=2I.
  \label{eq:double-layer-mass}
\end{equation}
Thus the raw positive layer has mass $2I$.  For
$\varphi\in C([0,2\pi])$, define
\begin{equation}
  \Phi_\gamma(\varphi)=\frac12\int_0^{2\pi}
  \varphi(t)D_B(t)\,dt.
  \label{eq:Phi-definition}
\end{equation}
The factor $1/2$ normalizes this mass.  The positivity of $D_B$
and~\eqref{eq:double-layer-mass} make $\Phi_\gamma$ positive and
unital, hence $*$-preserving.  It is bounded because
\[
  \lVert\Phi_\gamma(\varphi)\rVert
  \leq\frac12\lVert\varphi\rVert_\infty
  \int_0^{2\pi}\lVert D_B(t)\rVert\,dt.
\]
\end{proof}

\begin{proposition}[Radial Cayley completion]
\label{prop:cayley-completion}
Let $\Omega$ be a bounded convex domain with a supported oriented
radial boundary, let $B\in\Mn$ have simple spectrum, and suppose
$W(B)\subset\Omega$.  Let $V$ be an open neighborhood of
$\overline\Omega$, and let $f$ be holomorphic on $V$ and satisfy
\begin{equation}
  \max_{\zeta\in\overline\Omega}\lvert f(\zeta)\rvert\leq1.
  \label{eq:f-bounded-by-one}
\end{equation}
Put $T=f(B)$.  Then there is an analytic $H\colon\D\to\Mn$ satisfying
\begin{equation}
  H(0)=I,\qquad \RePart H(w)\succeq0\quad(w\in\D),
  \label{eq:cayley-H-positive}
\end{equation}
and
\begin{equation}
  H(w)-(I-wT)^{-1}\in\alg(B^*)\qquad(w\in\D).
  \label{eq:cayley-H-completion}
\end{equation}
\end{proposition}

\begin{proof}
Choose
\[
  B=S\diag(\beta_1,\ldots,\beta_n)S^{-1}.
\]
Then
$T=S\diag(f(\beta_1),\ldots,f(\beta_n))S^{-1}$, and
\eqref{eq:f-bounded-by-one} gives
$\lvert f(\beta_i)\rvert\leq1$ for every $i$.  In particular,
$\sigma(T)\subset\overline\D$.

For the fixed boundary data, put
$\mathcal R_t=(\gamma(t)I-B)^{-1}$, and let $\Phi_\gamma$ be the bounded
map from Lemma~\ref{lem:double-layer-map}.  For $w\in\D$, define
\begin{equation}
  c_w(t)=\frac{1+w f(\gamma(t))}{1-w f(\gamma(t))}
  \label{eq:cayley-boundary-function}
\end{equation}
and put
\begin{equation}
  H(w)=\Phi_\gamma(c_w).
  \label{eq:H-from-Phi}
\end{equation}
The expansion
\begin{equation}
  c_w(t)=1+2\sum_{m=1}^{\infty}w^m f(\gamma(t))^m
  \label{eq:cayley-uniform-series}
\end{equation}
converges locally uniformly in $C([0,2\pi])$.  The boundedness of
$\Phi_\gamma$ therefore makes $H$ analytic, and $H(0)=I$.  Furthermore,
\[
  \RePart c_w(t)
  =\frac{1-\lvert w f(\gamma(t))\rvert^2}
  {\lvert1-w f(\gamma(t))\rvert^2}\geq0.
\]
The positivity and $*$-preservation of $\Phi_\gamma$ give
$\RePart H(w)=\Phi_\gamma(\RePart c_w)\succeq0$.

For every $m\geq0$, the scalar Cauchy formula applied in the chosen
eigenbasis gives exactly
\begin{equation}
  \frac{1}{2\pi}\int_0^{2\pi}
  f(\gamma(t))^m\nu(t)\mathcal R_t s(t)\,dt=T^m.
  \label{eq:power-cauchy-identity}
\end{equation}
Indeed, this is the diagonal collection of the scalar Cauchy formulas
for $f^m$ at the points $\beta_i$; the neighborhood $V$ contains
the contour and its interior.  Transporting the uniformly convergent
series~\eqref{eq:cayley-uniform-series} through this first-layer
integral gives
\[
  I+2\sum_{m=1}^{\infty}w^mT^m
  =(I+wT)(I-wT)^{-1}.
\]
The matrix series converges because the powers of $T$ are bounded in
the displayed eigenbasis.

Define the companion matrix
\[
  C_f(w)=\frac{1}{2\pi}\int_0^{2\pi}
  \overline{c_w(t)}\,\nu(t)\mathcal R_t s(t)\,dt.
\]
Splitting the density~\eqref{eq:double-layer-density} into its two
layers now gives
\[
  2H(w)=(I+wT)(I-wT)^{-1}+C_f(w)^*.
\]
Each $\mathcal R_t$ belongs to $\alg(B)$, and this finite-dimensional
algebra is closed, so $C_f(w)\in\alg(B)$.  Since
$(I+wT)(I-wT)^{-1}=2(I-wT)^{-1}-I$, we conclude that
\begin{equation}
  H(w)-(I-wT)^{-1}
  =\frac12\bigl(C_f(w)^*-I\bigr)\in\alg(B^*).
  \label{eq:direct-completion}
\end{equation}
This proves~\eqref{eq:cayley-H-completion}.
\end{proof}

\begin{corollary}[Auxiliary sharp bound]
\label{cor:auxiliary-bound}
Under the hypotheses of Proposition~\ref{prop:cayley-completion},
$\lVert f(B)\rVert\leq2$.
\end{corollary}

\begin{proof}
Write
\[
  B=S\diag(\beta_1,\ldots,\beta_n)S^{-1}.
\]
Then
\[
  T=f(B)
  =S\diag\bigl(f(\beta_1),\ldots,f(\beta_n)\bigr)S^{-1}.
\]
Since $\sigma(B)\subset\Omega$, the bound on $f$ gives
$\lvert f(\beta_i)\rvert\leq1$ for every $i$.
Apply Theorem~\ref{thm:completion} with $\mathfrak m=2$ to these simultaneous
diagonalizations and to the function $H$ constructed in
Proposition~\ref{prop:cayley-completion}.  The values $f(\beta_i)$ may
repeat.
\end{proof}

\section{The finite-dimensional theorem}
\label{sec:passage}

Successive approximations pass from the auxiliary setting to arbitrary
matrices.

\begin{lemma}[Simple spectrum and numerical range]
\label{lem:simple-spectrum-density}
Matrices with $n$ distinct eigenvalues are dense in $\Mn$.  Moreover,
for $A,B\in\Mn$,
\begin{equation}
  W(B)\subset
  W(A)+\{z\in\C:\lvert z\rvert\leq\lVert B-A\rVert\}.
  \label{eq:numerical-range-Lipschitz}
\end{equation}
\end{lemma}

\begin{proof}
The discriminant of the characteristic polynomial is a nonzero
polynomial in the matrix entries, as is seen at
$\diag(1,\ldots,n)$.  The zero set of a nonzero complex polynomial has
empty interior, so its nonvanishing set is dense in
$\Mn$, proving the first assertion.  For a unit vector $x$,
\[
  \lvert x^*(B-A)x\rvert\leq\lVert B-A\rVert,
\]
which proves~\eqref{eq:numerical-range-Lipschitz}.
\end{proof}

\begin{lemma}[Parallel convex domains]
\label{lem:smooth-outer-approximation}
Let $K\subset\C$ be nonempty, compact, and convex.  For $\delta>0$, the
set
\[
  \Omega_\delta=\{z\in\C:\distop(z,K)<\delta\}
\]
is a bounded convex domain,
$\overline\Omega_\delta=K+\delta\overline\D$, and
$\Omega_\delta$ has a supported oriented radial boundary.
\end{lemma}

\begin{proof}
The set $\Omega_\delta$ is bounded and convex, and its closure is the
Minkowski sum $K+\delta\overline\D$.
Identify $\C$ with $\mathbb R^2$.  The metric projection
$\pi_K\colon\C\to K$ is unique and continuous.  Put
$d_K(z)=\distop(z,K)$.  For $z,h\in\C$, comparison first with $\pi_K(z)$
and then with $\pi_K(z+h)$ gives
\[
\begin{split}
  2\RePart\bigl(\overline{z+h-\pi_K(z+h)}\,h\bigr)-\lvert h\rvert^2
  &\leq d_K(z+h)^2-d_K(z)^2\\
  &\leq
  2\RePart\bigl(\overline{z-\pi_K(z)}\,h\bigr)+\lvert h\rvert^2.
\end{split}
\]
The continuity of $\pi_K$ now shows that $d_K^2$ is continuously
differentiable, with real gradient $2(z-\pi_K(z))$.  On the level set
$d_K(z)=\delta$ this gradient has norm $2\delta>0$.  Since
$\partial\Omega_\delta=\{d_K=\delta\}$, the
implicit-function theorem shows that $\partial\Omega_\delta$ is
continuously differentiable.

Choose $c\in K$.  The ball $c+\delta\D$ lies in $\Omega_\delta$, so
every ray from $c$ meets the boundary once and transversely.  Its radial
function $\rho$ is therefore positive, continuously differentiable,
and periodic.  With
\[
  \nu(t)=\frac{\gamma(t)-\pi_K(\gamma(t))}{\delta}
  \qquad(0\leq t\leq2\pi)
\]
as the outward unit normal along $\gamma$ and
$s(t)=\lvert\gamma'(t)\rvert$, convexity gives
the support inequality in~\eqref{eq:oriented-radial-boundary}; positive
orientation gives $\gamma'(t)=\iu\nu(t)s(t)$.  These data form the required
supported oriented radial boundary.
\end{proof}

\begin{proof}[Proof of Theorem~\ref{thm:main}]
We prove the holomorphic estimate~\eqref{eq:main-bound}; the polynomial
estimate~\eqref{eq:matrix-polynomial-bound} follows by taking $f=p$.
Put $K=W(A)$, and let $U$ be an open neighborhood of $K$ on which $f$
is holomorphic.  Choose $\varepsilon_0>0$ such that the closed
$\varepsilon_0$-neighborhood of $K$ lies in $U$.  For
$0<\varepsilon<\varepsilon_0$, let
\[
  \Omega_\varepsilon
  =\{z\in\C:\distop(z,K)<\varepsilon\}.
\]
Lemma~\ref{lem:smooth-outer-approximation} makes this a bounded convex
domain with a supported oriented radial boundary.  The open convex set
$V=\Omega_{\varepsilon_0}$ satisfies
$\overline\Omega_\varepsilon\subset V\subset U$.

By Lemma~\ref{lem:simple-spectrum-density}, choose simple-spectrum
matrices $B_k\to A$.  For fixed $\varepsilon$,
all sufficiently large $k$ satisfy
$\lVert B_k-A\rVert<\varepsilon/2$, and then
\eqref{eq:numerical-range-Lipschitz} gives
\[
  W(B_k)\subset K+(\varepsilon/2)\overline\D
  \subset\Omega_\varepsilon.
\]
Set
\[
  m_\varepsilon
  =\max_{z\in\overline\Omega_\varepsilon}\lvert f(z)\rvert.
\]
If $m_\varepsilon=0$, then $f$ vanishes on $\Omega_\varepsilon$, so
$f(A)=0$.  If $m_\varepsilon>0$, apply
Corollary~\ref{cor:auxiliary-bound} to $B_k$ and
$f/m_\varepsilon$ to obtain
\[
    \lVert f(B_k)\rVert\leq2m_\varepsilon.
\]
Continuity of the holomorphic functional calculus, obtained from a
fixed contour inside $U$, gives $f(B_k)\to f(A)$ and hence
the same upper bound for $f(A)$.  Thus, in either case,
\begin{equation}
  \lVert f(A)\rVert\leq2m_\varepsilon.
  \label{eq:epsilon-bound}
\end{equation}

Put $m_0=\max_{z\in K}\lvert f(z)\rvert$.  Uniform continuity of $f$
on the closed $\varepsilon_0$-neighborhood of $K$ gives a modulus of
continuity $\omega_f$ there.  Since every point of
$\overline\Omega_\varepsilon$ lies within $\varepsilon$ of $K$,
\[
  0\leq m_\varepsilon-m_0\leq\omega_f(\varepsilon).
\]
Letting $\varepsilon\downarrow0$ in~\eqref{eq:epsilon-bound} proves
\eqref{eq:main-bound}.

Finally, a rational function pole-free on the compact set $W(A)$ is
holomorphic on a neighborhood of $W(A)$.  Applying
\eqref{eq:main-bound} to such functions proves
\eqref{eq:spectral-set-definition} with $K=W(A)$ and $C=2$.
\end{proof}

\section{Consequences and scope}
\label{sec:consequences}

The finite-dimensional polynomial estimate passes to bounded
Hilbert-space operators by the standard compression argument
from~\cite{Crouzeix2007}; compare the survey~\cite{BickelEtAl2020}.
Polynomial approximation then gives the corresponding holomorphic and
spectral-set forms.

\begin{corollary}[Operator-level form]
\label{cor:operator-bound}
Let $\mathcal H\neq\{0\}$ be a complex Hilbert space, with inner
product linear in the first variable, and let $\mathcal B(\mathcal H)$
denote its algebra of bounded operators.  For
$A\in\mathcal B(\mathcal H)$, put
\[
  W(A)=\{\langle Ax,x\rangle:x\in\mathcal H,\ \lVert x\rVert=1\},
  \qquad K=\overline{W(A)}.
\]
Then, for every $p\in\C[z]$,
\begin{equation}
  \lVert p(A)\rVert
  \leq 2\sup_{z\in W(A)}\lvert p(z)\rvert.
  \label{eq:operator-polynomial-bound}
\end{equation}
More generally, if $f$ is holomorphic on a neighborhood of $K$, then
\begin{equation}
  \lVert f(A)\rVert
  \leq2\max_{z\in K}\lvert f(z)\rvert
  =2\sup_{z\in W(A)}\lvert f(z)\rvert.
  \label{eq:operator-holomorphic-bound}
\end{equation}
Consequently, $K$ is a $2$-spectral set for $A$.
\end{corollary}

\begin{proof}
We first prove~\eqref{eq:operator-polynomial-bound}.  The assertion is
immediate if $p=0$.  Otherwise, write $d=\deg p$.
For a unit vector $x\in\mathcal H$, let
\[
  \mathcal M_x=\operatorname{span}\{x,Ax,\ldots,A^d x\},
  \qquad
  A_x=\left.\Pi_x A\right|_{\mathcal M_x},
\]
where $\Pi_x$ is the orthogonal projection onto
$\mathcal M_x$.  Induction gives
\begin{equation}
  A_x^k x=A^k x \qquad(0\leq k\leq d),
  \label{eq:Krylov-agreement}
\end{equation}
because $A^{k+1}x\in\mathcal M_x$ whenever $k<d$.  Moreover,
$W(A_x)\subset W(A)$.  Theorem~\ref{thm:main}, applied on the
finite-dimensional space $\mathcal M_x$, therefore gives
\[
  \lVert p(A)x\rVert_{\mathcal H}
  =\lVert p(A_x)x\rVert_{\mathcal H}
  \leq\lVert p(A_x)\rVert
  \leq2\sup_{z\in W(A)}\lvert p(z)\rvert.
\]
Taking the supremum over unit vectors $x$ proves~%
\eqref{eq:operator-polynomial-bound}.

The set $K$ is compact and convex, and $\sigma(A)\subset K$.  Let $f$
be holomorphic on an open neighborhood $U$ of $K$.  Choose $\eta>0$
such that
\[
  L_\eta=\{z\in\C:\distop(z,K)\leq\eta\}\subset U.
\]
The compact convex set $L_\eta$ has connected complement, so polynomial
Runge approximation gives polynomials $p_j$ converging uniformly to
$f$ on $L_\eta$.  Let
\[
  \Gamma=\partial\{z\in\C:\distop(z,K)<\eta/2\},
\]
oriented positively.  Lemma~\ref{lem:smooth-outer-approximation}
shows that $\Gamma$ is a continuously differentiable contour enclosing
$\sigma(A)$.  The holomorphic functional calculus gives
\[
  p_j(A)-f(A)
  =\frac{1}{2\pi\iu}\int_\Gamma
  \bigl(p_j(z)-f(z)\bigr)(zI-A)^{-1}\,dz,
\]
and hence $p_j(A)\to f(A)$ in operator norm.  Since
\[
  \sup_{z\in W(A)}\lvert p_j(z)\rvert
  =\max_{z\in K}\lvert p_j(z)\rvert,
\]
passing to the limit in~\eqref{eq:operator-polynomial-bound} proves
\eqref{eq:operator-holomorphic-bound}.  Every rational function
pole-free on $K$ is holomorphic on a neighborhood of $K$, so the
spectral-set assertion follows.
\end{proof}

The finite-dimensional holomorphic estimate also controls polynomial
and rational approximation of matrix functions.  Let $A\in\Mn$, let
$f$ be holomorphic on a
neighborhood of $W(A)$, and let $s$ be either a polynomial or a
rational function pole-free on $W(A)$.  Then $f-s$ is holomorphic on
a neighborhood of $W(A)$, and Theorem~\ref{thm:main} gives
\begin{equation}
  \lVert f(A)-s(A)\rVert
  \leq2\max_{z\in W(A)}\lvert f(z)-s(z)\rvert.
  \label{eq:matrix-function-error}
\end{equation}
This is the sharp universal form of the numerical-range estimate that
underlies applications to GMRES and rational Krylov methods; see~%
\cite{ChoiGreenbaum2015,CrouzeixGreenbaum2019,ChenGreenbaumTrogdon2025}
for the role of spectral sets in those settings.

\section{Discussion: sharp consequences and open boundaries}
\label{sec:outlook}

The completion theorem separates a finite-dimensional algebraic
engine from a geometric positive layer.  We formulate the engine as a
reusable criterion and give detailed applications with matching upper
and lower constructions.  The closing subsection lists related open
sharpness questions.

\subsection{The positive-layer criterion and the abstract problem}

Let $\mathcal A$ be a uniform algebra on a compact space $K$, let
$\theta\colon\mathcal A\to\Mn$ be a continuous unital homomorphism,
and suppose that there is a simple-spectrum matrix $B$ such that
$\theta(\mathcal A)\subset\alg(B)$.  The following criterion is the
functional-calculus form of Theorem~\ref{thm:completion}.

\begin{proposition}[Positive-layer transfer]
\label{prop:positive-layer-transfer}
Suppose $2\leq\mathfrak m<\infty$ and that a bounded complex-linear map
$\mathcal P\colon\mathcal A\to\Mn$ satisfies
\begin{equation}
  \mathcal P(1)=\mathfrak m I,
  \qquad
  \RePart h\geq0\ \text{on }K
  \Longrightarrow\RePart\mathcal P(h)\succeq0,
  \qquad
  \mathcal P(h)-\theta(h)\in\alg(B^*)
  \label{eq:discussion-positive-layer}
\end{equation}
for every $h\in\mathcal A$.  Then
$\lVert\theta\rVert\leq\mathfrak m$.
\end{proposition}

\begin{proof}
Fix $f\in\mathcal A$ with $\lVert f\rVert_K\leq1$, put
$T=\theta(f)$, and, for $w\in\D$, define
\[
  c_w=(1+wf)(1-wf)^{-1},
  \qquad H(w)=\mathfrak m^{-1}\mathcal P(c_w).
\]
The diagonal entries of $T$ in the auxiliary basis are contractive
character values.  Moreover, $H$ is analytic, $H(0)=I$, and
$\RePart H(w)\succeq0$.  Since
\[
  \theta(c_w)=2(I-wT)^{-1}-I,
\]
condition~\eqref{eq:discussion-positive-layer} gives
\[
  H(w)-\frac2{\mathfrak m}(I-wT)^{-1}
  =\frac1{\mathfrak m}
   \bigl(\mathcal P(c_w)-\theta(c_w)-I\bigr)
  \in\alg(B^*).
\]
Theorem~\ref{thm:completion} yields
$\lVert\theta(f)\rVert\leq\mathfrak m$; taking the supremum over the
unit ball of $\mathcal A$ proves the assertion.
\end{proof}

The criterion asks for four verifiable items: positivity, total mass,
an adjoint-algebra defect, and an auxiliary simple-spectrum model,
together with approximation when needed.  A
complementary abstract approach based on extremal pairs and their
representing extremal measures is developed in
\cite{SchwenningerDeVries2024}.  One sharp abstract consequence of the
present criterion is the semisimple-range case of the problem posed in
\cite{ClouatreOstermannRansford2023}.

\begin{corollary}[Semisimple abstract Crouzeix problem]
\label{cor:abstract-semisimple-discussion}
Let $\alpha\colon\mathcal A\to\mathcal A$ be a unital antilinear map
and put
\[
  \theta_\alpha(h)=\frac12\bigl(
  \theta(h)+\theta(\alpha(h))^*\bigr).
\]
If $\lVert\theta_\alpha\rVert\leq1$ and the finite-dimensional
commutative algebra $\theta(\mathcal A)$ is semisimple, then
$\lVert\theta\rVert\leq2$.  The coefficient $2$ is best possible in
this semisimple subclass.
\end{corollary}

\begin{proof}
For each unit vector $x\in\C^n$, the scalar functional
$h\mapsto x^*\theta_\alpha(h)x$ is unital and contractive, hence extends
to a state on $C(K)$.  Thus $\theta_\alpha$ preserves nonnegative real
parts, and $\mathcal P=2\theta_\alpha$ is a positive layer of mass two
with
\[
  \mathcal P(h)-\theta(h)=\theta(\alpha(h))^*.
\]
Simultaneously diagonalize the semisimple range and choose, in the
same basis, an auxiliary matrix with distinct diagonal entries.
Proposition~\ref{prop:positive-layer-transfer} gives the upper bound.

For sharpness, take $0<\varepsilon,t<1$, put
\[
 T_\varepsilon=
 \begin{pmatrix}
  \varepsilon&2\sqrt{1-\varepsilon^2}\\0&-\varepsilon
 \end{pmatrix},
 \qquad
 T_{\varepsilon,t}=tT_\varepsilon,
 \qquad
 \theta_{\varepsilon,t}(h)=h(T_{\varepsilon,t}),
 \qquad
 \alpha(h)=\overline{h(0)}1
\]
on the disk algebra $\mathcal A=A(\D)$.  The elliptic range of
$T_\varepsilon$ is
contained in $\overline\D$, so $W(T_{\varepsilon,t})\subset\D$.
Apply Lemma~\ref{lem:double-layer-map} to the unit circle.  The Cauchy
calculation in the proof of Proposition~\ref{prop:cayley-completion}
gives
$2\Phi_\gamma(h\circ\gamma)=h(T_{\varepsilon,t})+h(0)I$.  Since
$\Phi_\gamma$ is unital and positive, the disk double-layer map
\[
 h\longmapsto\Phi_\gamma(h\circ\gamma)
 =\frac12\bigl(h(T_{\varepsilon,t})+h(0)I\bigr)
 =(\theta_{\varepsilon,t})_\alpha(h)
\]
is unital and contractive.  Since $T_{\varepsilon,t}$ has distinct
eigenvalues,
$\theta_{\varepsilon,t}(\mathcal A)=\alg(T_{\varepsilon,t})$ is
semisimple.  Finally,
\[
 \lVert\theta_{\varepsilon,t}\rVert
 \geq\lVert T_{\varepsilon,t}\rVert
 =t\bigl(1+\sqrt{1-\varepsilon^2}\bigr).
\]
Letting first $t\uparrow1$ and then $\varepsilon\downarrow0$ proves
sharpness.
\end{proof}

\subsection{A sharp higher-mass disk problem}

For $\varrho\geq1$, let $C_\varrho$ denote the class of bounded
operators $A\in\mathcal B(\mathcal H)$ for which there are a Hilbert
space $\mathcal K\supset\mathcal H$ and a unitary
$U\in\mathcal B(\mathcal K)$ such that
\[
  A^k=\varrho P_{\mathcal H}U^k|_{\mathcal H},
  \qquad k\geq1.
\]
Here $P_{\mathcal H}$ is the orthogonal projection of $\mathcal K$
onto $\mathcal H$.
For $\varrho\geq2$, the associated positive layer gives the following
exact constant.

\begin{theorem}[$C_\varrho$ disk calculus]
\label{thm:rho-contraction-discussion}
For each $\varrho\geq2$, the least universal constant for matrices
$A\in C_\varrho$ is $\varrho$:
\begin{equation}
  \lVert f(A)\rVert\leq\varrho
  \max_{\lvert z\rvert\leq1}\lvert f(z)\rvert
  \label{eq:rho-contraction-discussion}
\end{equation}
for every rational function $f$ pole-free on the closed disk.
The same sharp estimate on arbitrary Hilbert spaces is the classical
Okubo--Ando theorem~\cite{OkuboAndo1975}.
\end{theorem}

\begin{proof}
First suppose that $\sigma(A)\subset\D$.  On $\lvert\zeta\rvert=1$
put
\[
  P_A(\zeta)=\frac1\pi\RePart\bigl[
  \zeta(\zeta I-A)^{-1}\bigr].
\]
The Poisson-kernel characterization of $C_\varrho$ in
\cite{SchwenningerDeVriesRho2025} gives
$P_A(\zeta)\succeq(2-\varrho)(2\pi)^{-1}I$.  Hence
\[
  D_A(\zeta)=P_A(\zeta)-\frac{2-\varrho}{2\pi}I\succeq0.
\]
The disk double-layer identity gives, for the disk algebra,
\begin{equation}
  \int_{\lvert\zeta\rvert=1}D_A(\zeta)\,ds=\varrho I,
  \qquad
  \int_{\lvert\zeta\rvert=1}h(\zeta)D_A(\zeta)\,ds
  =h(A)+(\varrho-1)h(0)I.
  \label{eq:rho-positive-layer}
\end{equation}
The defect is scalar.  Proposition~\ref{prop:positive-layer-transfer}
therefore proves~\eqref{eq:rho-contraction-discussion} for
simple-spectrum matrices with spectrum in $\D$.

Now let $A\in C_\varrho\cap\Mn$ be arbitrary and define, for $z\in\D$,
\[
  Q_A(z)=\frac1\pi\RePart(I-zA)^{-1}
  -\frac{2-\varrho}{2\pi}I.
\]
The same characterization gives $Q_A(z)\succeq0$ and
$Q_A(0)=\varrho(2\pi)^{-1}I$.  Applying the scalar Harnack inequality
to every quadratic form gives, for $0<t<1$ and $\lvert z\rvert\leq t$,
\[
  Q_A(z)\succeq
  \frac{1-t}{1+t}\frac{\varrho}{2\pi}I.
\]
For $0\leq s\leq1$ and $\lvert\zeta\rvert=1$,
$D_{stA}(\zeta)=Q_A(st\overline\zeta)$, so the full density family for
$tA$ is uniformly positive.  Joint continuity in the matrix, $s$, and
$\zeta$ shows that every sufficiently small perturbation of $tA$ has
positive densities for $0\leq s\leq1$ and spectrum in $\D$, hence
belongs to $C_\varrho$.  Choose simple-spectrum matrices $A_j\to tA$
with this property.  The first part of the proof gives
\[
  \lVert f(A_j)\rVert\leq\varrho
  \max_{\lvert z\rvert\leq1}\lvert f(z)\rvert.
\]
Letting $A_j\to tA$ and then $t\uparrow1$ proves the upper bound.

For optimality, put
$J=\left(\begin{smallmatrix}0&1\\0&0\end{smallmatrix}\right)$ and
$A=\varrho J$.  The Sz.-Nagy--Foia\c{s} criterion recorded in
\cite{SchwenningerDeVriesRho2025} becomes
\[
 I-2(1-\varrho^{-1})\RePart(\overline zA)
 -(2\varrho^{-1}-1)\lvert z\rvert^2A^*A\succeq0.
\]
For $A=\varrho J$, the determinant of the displayed $2\times2$
matrix is $1-\lvert z\rvert^2$, and its lower-right entry is
$1+\varrho(\varrho-2)\lvert z\rvert^2$.  The matrix is therefore
positive definite for $z\in\D$.  Thus $A\in C_\varrho$, and the test
$f(z)=z$ has ratio
$\lVert A\rVert=\varrho$.  This witness also proves sharpness of the
Hilbert-space statement.
\end{proof}

\subsection{Sharp annular problems}

Fix $R>1$ and write
\[
  \mathbb A_R=\{z\in\C:R^{-1}<\lvert z\rvert<R\},
  \qquad
  \overline{\mathbb A}_R=
  \{z\in\C:R^{-1}\leq\lvert z\rvert\leq R\}.
\]

\begin{theorem}[Quantum annulus]
\label{thm:quantum-annulus-discussion}
If $A$ is an invertible Hilbert-space operator with
$\lVert A\rVert,\lVert A^{-1}\rVert\leq R$, then
\begin{equation}
  \lVert f(A)\rVert\leq2
  \max_{z\in\overline{\mathbb A}_R}\lvert f(z)\rvert
  \label{eq:quantum-annulus-discussion}
\end{equation}
for every rational $f$ pole-free on
$\overline{\mathbb A}_R$.  The constant $2$ is optimal for every
fixed $R$, already among finite matrices.
\end{theorem}

\begin{proof}
Orient the outer circle counterclockwise and the inner circle
clockwise, and parametrize each by oriented arclength $z=z(s)$.  For
a matrix $B$ satisfying
$\lVert B\rVert,\lVert B^{-1}\rVert<R$, define
\[
  \mu_B(z)=\frac1{2\pi\iu}\left[
  z'(s)(zI-B)^{-1}
  -\overline{z'(s)}(\overline zI-B^*)^{-1}\right],
  \qquad
  D_B(z)=\mu_B(z)-\frac1{2\pi\iu}\frac{z'(s)}zI.
\]
On the outer circle and on the inner circle, respectively, direct
multiplication gives
\begin{align}
  2\pi R D_B(z)
  &=(zI-B)^{-1}(R^2I-BB^*)
    (\overline zI-B^*)^{-1}\succeq0,\nonumber\\
  2\pi R^{-1}D_B(z)
  &=(zI-B)^{-1}(BB^*-R^{-2}I)
    (\overline zI-B^*)^{-1}\succeq0.
  \label{eq:annular-readiness-factorizations}
\end{align}
The two winding numbers cancel, while the two Cauchy terms each have
mass $I$.  Thus $D_B$ has mass $2I$, and the annular Cauchy identity
gives
\[
  \int h(z)D_B(z)\,ds-h(B)
  =(\mathcal C_R\overline h)(B)^*\in\alg(B^*),
\]
where
\[
 (\mathcal C_R\overline h)(\zeta)
 =\frac1{2\pi\iu}\int_{\partial\mathbb A_R}
   \frac{\overline{h(z)}}{z-\zeta}\,dz.
\]
To obtain the matrix upper bound, let $A\in\Mn$ satisfy
$\lVert A\rVert,\lVert A^{-1}\rVert\leq R$.  Choose $R'>R$ close
enough to $R$ that $f$ is pole-free on
$\overline{\mathbb A}_{R'}$, and choose simple-spectrum matrices
$B_j\to A$ with
$\lVert B_j\rVert,\lVert B_j^{-1}\rVert<R'$.  Let
$\mathcal A(\overline{\mathbb A}_{R'})$ be the algebra of functions continuous
on $\overline{\mathbb A}_{R'}$ and holomorphic on
$\mathbb A_{R'}$.  Write $D_{B_j}^{(R')}$ for the displayed density
with $R'$ in place of $R$, and set
$\theta_j(h)=h(B_j)$.  Then $\theta_j$ is a continuous unital
homomorphism into $\alg(B_j)$, and the positive map
\[
  \mathcal P_j\colon \mathcal A(\overline{\mathbb A}_{R'})
  \longrightarrow\Mn,\qquad
  \mathcal P_j(h)=
  \int_{\partial\mathbb A_{R'}}h(z)D_{B_j}^{(R')}(z)\,ds,
\]
has mass $2I$ and, for every
$h\in\mathcal A(\overline{\mathbb A}_{R'})$, satisfies
\[
  \mathcal P_j(h)-h(B_j)
  =(\mathcal C_{R'}\overline h)(B_j)^*\in\alg(B_j^*).
\]
Proposition~\ref{prop:positive-layer-transfer} yields
\[
  \lVert f(B_j)\rVert\leq2
  \max_{z\in\overline{\mathbb A}_{R'}}\lvert f(z)\rvert.
\]
Letting $j\to\infty$ and then $R'\downarrow R$ proves the matrix
estimate.

For an arbitrary Hilbert-space operator, the boundary-dilation theorem
of McCullough and Pascoe~\cite{McCulloughPascoe2023} provides a Hilbert
space $\mathcal K\supset\mathcal H$ and an invertible operator
$A_\partial\in\mathcal B(\mathcal K)$ such that
\begin{equation}
\begin{gathered}
 A^k=P_{\mathcal H}A_\partial^k|_{\mathcal H}
 \quad(k\in\mathbb Z),\\
 (R^2+R^{-2})I-A_\partial^*A_\partial
 -(A_\partial^*A_\partial)^{-1}=0.
\end{gathered}
 \label{eq:quantum-boundary-discussion}
\end{equation}
If $Z=A_\partial^*A_\partial$, then
$(Z-R^2I)(Z-R^{-2}I)=0$.  Hence
$Z=R^2E+R^{-2}(I-E)$ for a projection $E$.  Writing
$A_\partial=U_\partial Z^{1/2}$, let $u_0,p_0$ be the canonical unitary
and projection generators of $C(\mathbb T)*_{\C}\C^2$, where
$\mathbb T=\{z\in\C:\lvert z\rvert=1\}$, and put
\[
  j_R=u_0\bigl(Rp_0+R^{-1}(1-p_0)\bigr).
\]
The universal property gives a $*$-representation $\rho_\partial$ with
$\rho_\partial(u_0)=U_\partial$ and $\rho_\partial(p_0)=E$, and hence
$\rho_\partial(j_R)=A_\partial$.
Exel and Loring proved that this free product is residually
finite-dimensional~\cite{ExelLoring1992}.  Every finite-dimensional
$*$-representation $\pi$ sends $j_R$ and $j_R^{-1}$ to matrices of
norm at most $R$.  Thus, for every Laurent polynomial $q$,
\[
 \lVert q(A)\rVert
 \leq\lVert q(A_\partial)\rVert
 \leq\lVert q(j_R)\rVert
 =\sup_{\pi\ \mathrm{finite\ dimensional}}
   \lVert q(\pi(j_R))\rVert
 \leq2\max_{\overline{\mathbb A}_R}\lvert q\rvert.
\]
Choose $R''>R$ so that $f$ is pole-free on
$\overline{\mathbb A}_{R''}$, and choose Laurent polynomials $q_j$
converging uniformly to $f$ there.  The displayed estimate makes
$q_j(A)$ Cauchy, and the holomorphic functional calculus identifies
its limit with $f(A)$.  This proves
\eqref{eq:quantum-annulus-discussion}.

For sharpness, on $\C^{2n}$ let
\[
  W_ne_k=\omega_ke_{k+1\pmod{2n}},
  \qquad
  \omega_k=\begin{cases}
    R,&0\leq k<n,\\ R^{-1},&n\leq k<2n.
  \end{cases}
\]
Then $\lVert W_n\rVert=\lVert W_n^{-1}\rVert=R$.  For
$g_n(z)=R^{-n}(z^n+z^{-n})$,
\[
  \max_{\overline{\mathbb A}_R}\lvert g_n\rvert
  =1+R^{-2n},
  \qquad
  W_n^ne_0=W_n^{-n}e_0=R^ne_n,
\]
so $\lVert g_n(W_n)\rVert\geq2$.  The ratios tend to $2$.
This finite cyclic form complements the Hilbert-space lower-bound
construction in~\cite{Tsikalas2022}.
\end{proof}

The same exact constant persists for the centered double-layer class.
For an invertible matrix $A$ and $z,w\in\D$, use the notation
\[
 F_{R,0}^A(z,w)=
 2\RePart\left[(I-zA/R)^{-1}+(I-wA^{-1}/R)^{-1}\right]-2I.
\]
The class $\mathrm{DL}\mathbb A_R(0)$ consists of the matrices whose
spectrum lies in $\mathbb A_R$ and for which this kernel is
positive semidefinite on $\D^2$.

\begin{corollary}[Centered double-layer annulus]
\label{cor:centered-DL-discussion}
For every $R>1$, the least universal finite-dimensional constant for
the class $\mathrm{DL}\mathbb A_R(0)$ of Jury and Tsikalas is $2$.
Equivalently, every $A$ in that class and every rational $f$ pole-free
on $\overline{\mathbb A}_R$ satisfy
\[
 \lVert f(A)\rVert\leq2
 \max_{z\in\overline{\mathbb A}_R}\lvert f(z)\rvert.
\]
\end{corollary}

\begin{proof}
Let $A\in\mathrm{DL}\mathbb A_R(0)$ and fix such an $f$.  Choose
$R'>R$ sufficiently close to $R$ that $f$ remains pole-free on
$\overline{\mathbb A}_{R'}$.  The defining kernel satisfies
\[
 F_{R',0}^A(z,w)=
 F_{R,0}^A\bigl((R/R')z,(R/R')w\bigr),
 \qquad z,w\in\D.
\]
For each unit vector $x$, the function
\[
 (z,w)\longmapsto
 x^*F_{R,0}^A\bigl((R/R')z,(R/R')w\bigr)x
\]
is nonnegative and pluriharmonic on a neighborhood of
$\overline\D^2$, and its value at $(0,0)$ is $2$.  The strong minimum
principle makes it positive throughout $\overline\D^2$.  Compactness
of the unit sphere times $\overline\D^2$ yields a uniform positive
lower bound.  Hence all
sufficiently small simple-spectrum perturbations $A_j$ of $A$ belong
to $\mathrm{DL}\mathbb A_{R'}(0)$ and have spectrum in the open
annulus.  For each $A_j$, the normalized annular double-layer map
constructed by Jury and Tsikalas~\cite{JuryTsikalas2024} is unital and
contractive.  Its
unnormalized form is a positive layer of mass two, and the
double-layer identity leaves an adjoint Cauchy companion.
Proposition~\ref{prop:positive-layer-transfer} gives
\[
 \lVert f(A_j)\rVert\leq2
 \max_{z\in\overline{\mathbb A}_{R'}}\lvert f(z)\rvert.
\]
Let $A_j\to A$ and then $R'\downarrow R$ to obtain the upper bound.

Every quantum-annulus matrix whose spectrum lies in $\mathbb A_R$
belongs to $\mathrm{DL}\mathbb A_R(0)$.  Indeed, for a contraction $X$ and
$z\in\D$,
\[
  2\RePart(I-zX)^{-1}-I
  =(I-\overline zX^*)^{-1}
   (I-\lvert z\rvert^2X^*X)(I-zX)^{-1}\succeq0;
\]
apply this identity to $X=A/R$ and $X=A^{-1}/R$ and add.  The matrices
$W_n$ above provide the matching lower bound.
\end{proof}

For $c\in\C$ and $0<r<R$, let $A$ be a bounded Hilbert-space operator
such that $A-cI$ is invertible and
\[
  \lVert A-cI\rVert\leq R,
  \qquad
  \lVert(A-cI)^{-1}\rVert\leq r^{-1}.
\]
Then
\[
  \lVert f(A)\rVert\leq
  2\max_{r\leq\lvert z-c\rvert\leq R}\lvert f(z)\rvert
\]
for every rational function $f$ pole-free on the displayed annulus,
and the factor $2$ is optimal.  Apply
Theorem~\ref{thm:quantum-annulus-discussion} to
$(A-cI)/\sqrt{rR}$ with annular parameter $\sqrt{R/r}$, and transport
the weighted-shift lower sequence.

\subsection{Sharp geometric transfers}

The numerical-range theorem also solves the scaled $q$-range problem
by a rank-one similarity extraction.  We include the argument because
it is the step that determines the constant.

For $n\geq2$, $A\in\Mn$, and $0<\lvert q\rvert\leq1$, put
\[
  W_q(A)=\{y^*Ax:\lVert x\rVert_2=\lVert y\rVert_2=1, y^*x=q\},
  \qquad \Omega_q(A)=q^{-1}W_q(A).
\]
Let $r=\lvert q\rvert$ and
$\varkappa=(1+\sqrt{1-r^2})/r$.  A phase change in $y$ gives
$\Omega_q(A)=\Omega_r(A)$.  For a unit vector $x$, put
\[
 d_A(x)=\left(\lVert Ax\rVert_2^2-
 \lvert x^*Ax\rvert^2\right)^{1/2}.
\]
Tsing's disk formula~\cite{Tsing1984} is
\begin{equation}
 \Omega_r(A)=\bigcup_{\lVert x\rVert_2=1}
 \left\{z:\lvert z-x^*Ax\rvert\leq
 \frac{\sqrt{1-r^2}}r d_A(x)\right\}.
 \label{eq:q-disk-formula-discussion}
\end{equation}
It implies both
$W(A)\subset\Omega_r(A)$ and
$\Omega_s(A)\subset\Omega_r(A)$ for $r\leq s\leq1$.
For a unit vector $v$, let $P_v$ be the projection onto $\C v$ and
set $S_v=I+(\varkappa-1)P_v$.

\begin{lemma}[Rank-one stretch and extraction]
\label{lem:q-transfer-discussion}
For every unit $v$,
\begin{equation}
  W(S_vAS_v^{-1})\subset\Omega_r(A).
  \label{eq:q-stretch-discussion}
\end{equation}
If $X\in\Mn$, $\sigma(X)\subset\overline\D$, and
\begin{equation}
 \lVert S_v\varphi_{\omega,a}(X)S_v^{-1}\rVert\leq K
 \label{eq:q-extraction-hypothesis}
\end{equation}
for every unit $v$, $\lvert\omega\rvert=1$, $0\leq a<1$, where
$\varphi_{\omega,a}(z)=(\omega z-a)/(1-a\omega z)$, then
\begin{equation}
  \lVert X\rVert\leq\max\{1,K/\varkappa\}.
  \label{eq:q-extraction-conclusion}
\end{equation}
\end{lemma}

\begin{proof}
For a unit $u$, put
$\ell=\lVert S_v^{-1}u\rVert_2$ and
$m=\lVert S_vu\rVert_2$.  Directly from the two eigenvalues
$1,\varkappa$ of $S_v$,
\[
  1\leq\ell m\leq\tfrac12(\varkappa+\varkappa^{-1})=r^{-1}.
\]
With $x=S_v^{-1}u/\ell$, $y=S_vu/m$, and
$s=(\ell m)^{-1}$, one has $r\leq s\leq1$, $y^*x=s$, and
\[
  u^*S_vAS_v^{-1}u=s^{-1}y^*Ax\in\Omega_s(A)
  \subset\Omega_r(A).
\]
This proves~\eqref{eq:q-stretch-discussion}.

For the extraction, put $t=\lVert X\rVert$ and assume $t>1$.
Choose unit vectors $x,y$ with $Xx=ty$, and choose $\omega$ so that
$c=x^*(\omega y)=\lvert x^*y\rvert\in[0,1)$.  The strict inequality
follows because $c=1$ would give an eigenvalue of modulus $t>1$.
There is an
$a\in[0,1)$ satisfying
\begin{equation}
  tc(1+a^2)=a(1+t^2);
  \label{eq:q-orthogonality}
\end{equation}
take $a=0$ when $c=0$, and otherwise use the sign change between
$a=0$ and $a=1$.  Set
\[
  u=(I-a\omega X)x,
  \qquad w=(\omega X-aI)x.
\]
Since $I-a\omega X$ is invertible, $u\neq0$; moreover,
$\lVert\omega Xx\rVert_2=t>a$, so $w\neq0$.  Then
$\varphi_{\omega,a}(X)u=w$ and
\eqref{eq:q-orthogonality} gives $u\perp w$.  It also gives
\[
  \lVert w\rVert_2^2-t^2\lVert u\rVert_2^2
  =atc(t^2-1)(1-a^2)\geq0.
\]
Taking $v=w/\lVert w\rVert_2$, the stretch fixes $u$ and multiplies
$w$ by $\varkappa$.  Testing~\eqref{eq:q-extraction-hypothesis} on
$u$ yields $K\geq\varkappa t$, proving
\eqref{eq:q-extraction-conclusion}.
\end{proof}

\begin{theorem}[Scaled $q$-numerical ranges]
\label{thm:q-range-discussion}
For each fixed $q$ with $0<\lvert q\rvert\leq1$, the least constant
valid uniformly over all
$n\geq2$ and $A\in\Mn$ is
\begin{equation}
  C_q=\max\left\{1,
  \frac{2\lvert q\rvert}
  {1+\sqrt{1-\lvert q\rvert^2}}\right\}.
  \label{eq:q-sharp-discussion}
\end{equation}
More precisely, for every $n\geq2$, every $A\in\Mn$, and every
rational function $f$ pole-free on $\Omega_q(A)$,
\[
 \lVert f(A)\rVert\leq C_q
 \max_{z\in\Omega_q(A)}\lvert f(z)\rvert.
\]
Thus the polynomial conjecture proposed in
\cite{OLoughlinRani2026} holds, together with its rational spectral-set
form.
\end{theorem}

\begin{proof}
Let $M=\max_{\Omega_r(A)}\lvert f\rvert$, fix $\varepsilon>0$, and
put $X=f(A)/(M+\varepsilon)$.  For every unit $v$,
Lemma~\ref{lem:q-transfer-discussion} puts
$W(S_vAS_v^{-1})$ in $\Omega_r(A)$.  Apply
Theorem~\ref{thm:main} to the composition
\[
  z\longmapsto
  \varphi_{\omega,a}\bigl(f(z)/(M+\varepsilon)\bigr).
\]
Similarity covariance gives~\eqref{eq:q-extraction-hypothesis} with
$K=2$; spectral mapping gives $\sigma(X)\subset\D$.  The extraction
conclusion and
$2/\varkappa=2r/(1+\sqrt{1-r^2})$, followed by
$\varepsilon\downarrow0$, prove the upper bound.

Constant functions force $C_q\geq1$.  For the second branch, use
$J=\left(\begin{smallmatrix}0&1\\0&0\end{smallmatrix}\right)$.
Tsing's disk formula gives
\[
  \max_{z\in\Omega_r(J)}\lvert z\rvert=\frac{\varkappa}{2}.
\]
For a unit vector $x=(\xi,\zeta)$ put
$a_x=\lvert\xi\rvert$ and $b_x=\lvert\zeta\rvert$.  Then
$d_J(x)=b_x^2$, so every point in the disk indexed by $x$ has modulus
at most
\[
 a_xb_x+\frac{\varkappa^2-1}{2\varkappa}b_x^2
 \leq\frac{\varkappa}{2},
 \qquad
 \frac{\varkappa}{2}-a_xb_x
 -\frac{\varkappa^2-1}{2\varkappa}b_x^2
 =\frac{(\varkappa a_x-b_x)^2}{2\varkappa}.
\]
Equality is attained at
$a_x=(1+\varkappa^2)^{-1/2}$ and
$b_x=\varkappa(1+\varkappa^2)^{-1/2}$.  The identity polynomial has
ratio $2/\varkappa$, proving sharpness of~\eqref{eq:q-sharp-discussion}.
\end{proof}

Two further exact consequences concern fixed ellipses and the
Douglas--Paulsen norm.  For $0\leq\delta<1$, define
\[
  K_\delta=\left\{x+\iu y:
  \frac{x^2}{(1+\delta)^2}
  +\frac{y^2}{(1-\delta)^2}\leq1\right\}.
\]

\begin{theorem}[Fixed ellipses]
\label{thm:ellipse-discussion}
For every fixed $\delta\in[0,1)$, every bounded Hilbert-space operator
$T$ with $\overline{W(T)}\subset K_\delta$, and every rational function
$f$ pole-free on $K_\delta$,
\[
  \lVert f(T)\rVert\leq2
  \max_{z\in K_\delta}\lvert f(z)\rvert.
\]
The least constant uniform in $T$ and $f$ is $2$.  The same conclusion
holds for every nondegenerate ellipse after an affine change of
variables.
\end{theorem}

\begin{proof}
The upper bound follows from Corollary~\ref{cor:operator-bound}.  For
$0<\delta<1$, let $X_n$ be the cyclic weighted shift
on $\C^{2n}$ with $n$ consecutive weights $1$ followed by $n$ weights
$\delta$.  Then
\[
  X_n^{2n}=\delta^nI,
  \qquad \lVert X_n\rVert\leq1,
  \qquad \lVert X_n^{-1}\rVert\leq\delta^{-1}.
\]
Put $T_n=X_n+\delta X_n^{-1}$.  A numerical-range formula of Agler,
Lykova, and Young~\cite{AglerLykovaYoung2024} gives
$W(T_n)\subset K_\delta$.  Define
\[
  \mathcal F_0=2,\quad \mathcal F_1(z)=z,\quad
  \mathcal F_k(z)=z\mathcal F_{k-1}(z)
  -\delta\mathcal F_{k-2}(z)\quad(k\geq2).
\]
Induction gives
$\mathcal F_k(z+\delta z^{-1})=z^k+\delta^kz^{-k}$.
Since
$\partial K_\delta=\{e^{\iu t}+\delta e^{-\iu t}:0\leq t\leq2\pi\}$,
the maximum-modulus principle gives
\[
  \max_{K_\delta}\lvert\mathcal F_n\rvert=1+\delta^n,
  \qquad
  \mathcal F_n(T_n)=2X_n^n,
  \qquad \lVert X_n^n\rVert=1.
\]
Thus the ratios are $2/(1+\delta^n)$ and tend to $2$.  For
$\delta=0$, use $T=2J$ and the identity polynomial, where $J$ is
the matrix in~\eqref{eq:sharp-example}.
\end{proof}

Let $\mathcal R_\delta=\{z:\delta<\lvert z\rvert<1\}$, and for
$\varphi\in H^\infty(\mathcal R_\delta)$ let
$\lVert\varphi\rVert_{\rm dp}$ be the supremum of
$\lVert\varphi(X)\rVert$ over invertible $X$ with
$\lVert X\rVert\leq1$ and
$\lVert X^{-1}\rVert\leq\delta^{-1}$ and
$\sigma(X)\subset\mathcal R_\delta$, as in
\cite{AglerLykovaYoung2025}.

\begin{corollary}[Douglas--Paulsen norm]
\label{cor:DP-discussion}
For every $0<\delta<1$ and every
$\varphi\in H^\infty(\mathcal R_\delta)$,
\[
  \lVert\varphi\rVert_\infty
  \leq\lVert\varphi\rVert_{\rm dp}
  \leq2\lVert\varphi\rVert_\infty,
\]
and the factor $2$ is optimal.
\end{corollary}

\begin{proof}
Scalar operators give the first inequality.  For the second, put
$Y=\delta^{-1/2}X$ and apply
Theorem~\ref{thm:quantum-annulus-discussion} to
$f(z)=\varphi(\delta^{1/2}z)$ when $\varphi$ is holomorphic past the
closed annulus.  In general, write
$\varphi(z)=\sum_{k\in\mathbb Z}a_kz^k$ and use the angular Abel means
\[
 \varphi_s(z)=\sum_{k\in\mathbb Z}s^{\lvert k\rvert}a_kz^k,
 \qquad 0<s<1.
\]
They are Poisson averages of rotations of $\varphi$, so
$\lVert\varphi_s\rVert_\infty\leq\lVert\varphi\rVert_\infty$; they
are holomorphic on $s\delta<\lvert z\rvert<s^{-1}$.  Uniform Laurent
approximation therefore permits application of
Theorem~\ref{thm:quantum-annulus-discussion}, and the holomorphic
functional calculus gives $\varphi_s(X)\to\varphi(X)$ as
$s\uparrow1$.  This proves the upper bound for all of
$H^\infty(\mathcal R_\delta)$.  For the shifts $X_n$ above,
$X_n^{2n}=\delta^nI$ gives
$\sigma(X_n)\subset\{z:\lvert z\rvert=\sqrt\delta\}
\subset\mathcal R_\delta$.  The functions
$\chi_n(z)=z^n+\delta^nz^{-n}$ satisfy
$\lVert\chi_n\rVert_\infty=1+\delta^n$ and
$\lVert\chi_n(X_n)\rVert=2$, proving sharpness.
\end{proof}

Finally, a spherical disk is a closed disk, a closed half-plane, or the
closed exterior of a disk in the Riemann sphere.  Such a set $D$ is a
one-spectral set for $A$ if $\sigma(A)\subset D$ and
\[
  \lVert f(A)\rVert\leq\sup_{z\in D}\lvert f(z)\rvert
\]
for every rational function $f$ with poles outside $D$.  Let $A$ be a
bounded Hilbert-space operator and let $D_1,D_2$ be spherical disks,
each a one-spectral set for $A$; the general framework is developed in
\cite{BadeaBeckermannCrouzeix2009}.
The M\"obius reductions of Crouzeix and Greenbaum
\cite{CrouzeixGreenbaum2019} send the disjoint-boundary case to a
bounded-disk/exterior-disk pair and the crossing case to a
disk--half-plane pair.

\begin{corollary}[Two spherical disks]
\label{cor:two-disks-discussion}
The least constant valid uniformly for all such pairs is $2$:
\[
  \lVert f(A)\rVert\leq2
  \sup_{z\in D_1\cap D_2}\lvert f(z)\rvert.
\]
Here $f$ is any rational function pole-free on $D_1\cap D_2$.  If the
boundary circles are disjoint and $D_1\cap D_2$ is doubly connected,
the optimal constant for that fixed pair is $2$.
\end{corollary}

\begin{proof}
M\"obius covariance reduces the doubly connected case to a round
annulus, already handled above.  If the circles
cross, a M\"obius map with pole outside $D_1\cap D_2$ converts the
intersection into a disk--half-plane intersection containing the
numerical range of the transformed operator: the disk and half-plane
one-spectral hypotheses give this containment through their affine
and Cayley test functions.  The intersection is convex, so
Corollary~\ref{cor:operator-bound} applies.  Tangency follows by
enlargement and a limit, and a redundant constraint has constant one.
These cases give the upper bound.  Transported cyclic annular shifts
show that every doubly connected case has lower constant two, and
hence that the
uniform two-disk constant is exactly two.
\end{proof}

\subsection{Sharp formal transfers}

Faithful representations give two additional exact consequences.  If
$\mathfrak C$ is a unital $C^*$-algebra, define the algebraic numerical
range by
\[
  V_{\mathfrak C}(a)=
  \{\phi(a):\phi\in\mathfrak C^\vee,
  \ \lVert\phi\rVert=\phi(1)=1\}.
\]
$\mathfrak C^\vee$ denotes the continuous dual of $\mathfrak C$.
Then, for every
rational function $f$ pole-free on $V_{\mathfrak C}(a)$,
\begin{equation}
  \lVert f(a)\rVert\leq
  2\max_{z\in V_{\mathfrak C}(a)}\lvert f(z)\rvert.
  \label{eq:Cstar-discussion}
\end{equation}
A faithful universal representation identifies
$V_{\mathfrak C}(a)$ with the closure of the represented numerical
range, so Corollary~\ref{cor:operator-bound} proves the estimate.
The matrix algebra example $a=2J$ proves that the universal constant
over all unital $C^*$-algebras is exactly two.  Now let $\mathcal H$
be infinite-dimensional, let $\mathcal K(\mathcal H)$ denote the
compact operators, let $[A]$ be the coset of $A$ in the Calkin
algebra, and define
$W_{\mathrm e}(A)=V_{\mathcal B(\mathcal H)/\mathcal K(\mathcal H)}([A])$.
Applying the same argument to this quotient gives, for every rational
$f$ pole-free on $W_{\mathrm e}(A)$,
\[
  \lVert f([A])\rVert\leq
  2\max_{z\in W_{\mathrm e}(A)}\lvert f(z)\rvert.
\]
If $f(A)$ is defined, the left side is its essential norm.  Infinite
amplification of $2J$ proves sharpness; the same equality of universal
constants is established in~\cite{BlazhkoEtAl2025}.

Operator radii preserve all of the sharp scalar information above.
For $\upsilon\geq1$, put
\[
 w_\upsilon(T)=\inf\{s>0:s^{-1}T\in C_\upsilon\}.
\]
Let $\mathcal U$ be a unital Banach algebra satisfying
\[
  \lVert p(a)\rVert\leq
  \max_{\lvert z\rvert\leq1}\lvert p(z)\rvert
  \qquad(a\in\mathcal U,\ \lVert a\rVert\leq1,\ p\in\C[z]),
\]
and let $\Xi\colon\mathcal U\to\mathcal B(\mathcal H)$ be a bounded
unital homomorphism.  Put
\[
 \lVert\Xi\rVert_{w_\upsilon}
 =\sup_{\lVert h\rVert\leq1}w_\upsilon(\Xi(h)).
\]
Badea, Crouzeix, and Klaja prove the exact identity
\cite{BadeaCrouzeixKlaja2018}
$\lVert\Xi\rVert_{w_\upsilon}=H_\upsilon(\lVert\Xi\rVert)$, where, for
$s\geq1$,
\begin{equation}
  H_\upsilon(s)=
  \frac{s^2+1+
  \sqrt{(s^2+1)^2-4\upsilon(2-\upsilon)s^2}}
  {2\upsilon s};
  \label{eq:operator-radius-conversion-discussion}
\end{equation}
Each relevant functional-calculus domain satisfies the displayed von
Neumann inequality: the compact cases use uniform algebras, and the
Douglas--Paulsen calculus uses $H^\infty(\mathcal R_\delta)$.  Consequently
every sharp constant-two family above has exact
$w_\upsilon$-constant
\begin{equation}
  H_\upsilon(2)=
  \frac{5+\sqrt{25-16\upsilon(2-\upsilon)}}{4\upsilon},
  \qquad H_2(2)=\frac54,
  \label{eq:operator-radius-discussion}
\end{equation}
and the scaled $q$-range family has exact constant $H_\upsilon(C_q)$.
Continuity and strict monotonicity of $H_\upsilon$ transfer each lower
sequence and prove equality.  The $C_\varrho$ disk family similarly
gives the exact two-parameter constant $H_\upsilon(\varrho)$ for
$\varrho\geq2$.

\subsection{Related open problems}

Several nearby sharp-constant questions remain open.  The present mass
criterion gives the finite-dimensional upper bounds $2+c$ for the
noncentered classes $\mathrm{DL}\mathbb A_R(c)$, $c>0$, defined in
\cite{JuryTsikalas2024}, $N+1$ for a disk with $N\geq1$
pairwise disjoint circular holes lying strictly inside the outer disk
under the Badea--Beckermann--Crouzeix norm and inverse-resolvent
constraints, and $4$ under the numerical-annulus
conditions $w(A),w(A^{-1})\leq R$, where
$w(T)=\sup\{\lvert z\rvert:z\in W(T)\}$.  Determining matching lower
bounds is open.  Definitions and prior bounds appear in
\cite{BadeaBeckermannCrouzeix2009,
TsikalasMultiplyConnected2026,CrouzeixAnnuli2025}.
Further open questions concern the optimal constant for an individual
crossing lens, optimal fixed-domain constants for unbounded convex
numerical ranges, for which aperture-dependent upper bounds are given
in~\cite{CrouzeixUnbounded2025}, polyannuli and intersections of
three or more spherical disks, infinite-dimensional reduction for the
full double-layer classes, the complete matrix-valued numerical-range
problem~\cite{DavidsonPaulsenWoerdeman2018}, related abstract complete
reductions~\cite{HartzMcCarthy2026}, and the classification of
extremizers suggested by the multiplicity phenomena in~\cite{Li2020}.

\section*{AI-assistance disclosure}

OpenAI ChatGPT contributed the idea of sampling the matrix Herglotz
kernel at scaled conjugate eigenvalues, with scale $1/2$ in the
mass-two application, and adding the origin sample that cancels the
nonconstant adjoint-algebra correction's contribution.  This appears
in the proof of Theorem~\ref{thm:completion}, especially~%
\eqref{eq:eigenvalue-samples}--\eqref{eq:theta-contribution}.  It also
assisted with exposition, bibliography, and typesetting.  The human
author assumes full responsibility for the correctness of all
mathematical statements and for the integrity and accuracy of all
citations.

\section*{Acknowledgments}

The author thanks Elijah Winners for suggesting the auxiliary-basis
simplification used in the proof of
Theorem~\ref{thm:completion}.

\reftitle{References}
\bibliography{the_numerical_range_is_a_2_spectral_set}

\end{document}
